% =====================================================================
%  RESUMO EM LÍNGUA ESTRANGEIRA -- ABSTRACT (obrigatório) -- Manual, 2.1.8.2
%  É a TRADUÇÃO FIEL do resumo em português, incluindo as palavras-chave
%  (Keywords). Mesmas regras de formatação do resumo.
% =====================================================================
\begin{resumo}[Abstract]
\begin{otherlanguage*}{english}
\SingleSpacing
\noindent Dropout is a recurring problem in martial arts and combat sports
gyms, and many students quit before perceiving the physical and mental
benefits of training. This work presents RoundUp, a gamified mobile
application that aims to increase student engagement and retention and to
support instructors and gyms in managing their classes. The project was
conducted as an ideation process: thirteen feature ideas were converted
into requirements, which guided the development of the application and,
at the same time, were evaluated through a questionnaire answered by current
and former practitioners. To broaden the coverage of user profiles,
including the instructor profile, synthetic personas and a response
generator developed with the support of artificial intelligence were
employed: from 12 personas, 65 responses were produced, which made it
possible to exercise the entire questionnaire and the analysis procedure,
and were kept separate from the real data. The prototype was developed with React Native and Firebase, has
three user profiles (student, instructor and gym) and implements 29 of the 33 functional requirements, including attendance check-in validated by
the instructor, experience points, levels, training streaks, badges,
leaderboard, missions, belt promotion records, attendance reports and
consent in compliance with the Brazilian data protection law (LGPD). The
business rules that protect the gamification were implemented in the
database security rules and covered by automated tests. In the real
evaluation, the features that make progress visible, such as the progress
bar toward the next belt and the attendance history, received the highest
scores; reminders and automatic messages, the lowest; and the comments
pointed to excessive notifications and data entry as reasons to abandon
the app, which led to making notifications configurable and the
leaderboard optional. It is concluded that gamification centered on
visible progress is the most promising direction for the application,
that part of the dropout has external causes that an app cannot solve,
and that the next step is to validate the management requirements with
instructors in the field.

\vspace{\onelineskip}
\noindent\textbf{Keywords}: gamification; martial arts; retention; requirements engineering; personas.
\end{otherlanguage*}
\end{resumo}
